\documentclass[10.5pt]{config}
% 本实验报告采用 署名-相同方式共享 4.0 国际许可协议 (CC BY-SA 4.0)
\usepackage{multirow}
\usepackage{wrapfig}
\usepackage{setspace}
\usepackage{titlesec}
\usepackage{ctex}
\usepackage{graphicx} % 用于插入图片
\usepackage{float}    % 用于控制图片位置
\usepackage{makecell}

\titleformat{\subsubsection}[block]  % 设置为块级标题（自动换行）
{\normalfont\bfseries}{\thesubsubsection}{1em}{} 
\titlespacing*{\subsubsection}{1.5em}{3.25ex}{1.5ex}  % 调整缩进：左缩进2em
\setstretch{1.5}

% --- 请在这里填写您的个人信息 ---
\major{机械工程} % 专业
\name{韩颂义}
\partner{} % 签名栏
\stuid{3240104198}
\date{4.2}
\lab{东3-208}
\grades{}
\course{电工电子学实验}
\instructor{季瑞松} % 指导老师
\exptype{基础实验} % 实验类型
\expname{单相交流电特性及功率因数提高} % 实验名字


\begin{document}

\makeheader % 首页头部

\section{实验目的和要求}

\begin{enumerate}[label=\arabic*., leftmargin=2em]
    \item 了解电感性电路提高功率因数的方法和意义。
    \item 学会使用交流仪表（电压表、电流表、功率表）。
    \item 掌握用交流仪表测量交流电路电压、电流和功率的方法。
\end{enumerate}

\section{实验内容和原理}


\subsection{实验原理}

交流电路中常用的无源电路器件有电阻、电感和电容，利用交流电压表、交流电流表、功率表测量有关的电压、电流和功率，可以算得被测电路的有关参数。交流电路中的负载通常有电阻性负载，如白炽灯、电阻加热器等，也有电感性负载，如电动机、变压器、电磁线圈等。电磁线圈通常带有铁芯，在交流工况下铁芯会产生一定的损耗，因此一个具有铁芯的线圈在交流工况下其等效电阻是不等于直流工况下的等效电阻的。

\subsubsection*{1. 线圈参数的测量}

将一个 $R$ 値电阻与线圈串联后接在工频交流电源上（开关 S 断开），用交流电压表和交流电流表分别测量有效値 $U$、$U_R$、$U_L$、$I$（$=I_{RL}$）。以电源电压相量 $\dot{U}$ 为参考相量，由相量关系可得：
\begin{equation*}
    \cos\varphi_{RL} = \frac{U^2 + U_R^2 - U_L^2}{2U_R U}, \quad
    \sin\varphi_{RL} = \sqrt{1 - \cos^2\varphi_{RL}}
\end{equation*}
\begin{equation*}
    X_L = \frac{U\sin\varphi_{RL}}{I_{RL}} \quad \text{或} \quad
    L = \frac{U\sin\varphi_{RL}}{\omega I_{RL}}, \quad
    R_L = \frac{U\cos\varphi_{RL} - U_R}{I_{RL}}
\end{equation*}

\subsubsection*{2. 并联电容提高功率因数}

工业上大量的设备为电感性负载，由于电感性负载存在较大的感抗，因而电路的功率因数较低。对于电感性负载，可以采用并联电容器的方法来提高电路的功率因数，从而提高电源设备的利用率，降低输电线路中的损耗。

当开关 S 闭合时，在外加正弦交流电压 $\dot{U}$ 的作用下，各支路电流的关系为
\begin{equation*}
    \dot{I} = \dot{I}_{RL} + \dot{I}_C
\end{equation*}
并联容量适当的电容 $C$ 后，电容电流 $\dot{I}_C$ 补偿了电流 $\dot{I}_{RL}$ 的部分无功分量，使得电路总电流 $\dot{I}$ 的数値减小，$\dot{U}$ 与 $\dot{I}$ 的相位差 $\varphi$ 减小，功率因数提高。但若并联电容的容量过大，总电流 $\dot{I}$ 又会增大，功率因数下降，所以与电感性电路并联的电容量要适当。

\subsubsection*{3. 功率表的使用}

交流电路所消耗的有功功率可以用功率表（也称瓦特表）来测量。功率表内部包含电流测量部件与电压测量部件。电流测量部件工作时与被测电路相串联，电压测量部件工作时与被测电路相并联。测量负载功率时，电压测量端口的“$*$”端和电流测量端口的“$*$”端连接在一起后接到电源端，读数为正表示被测电路吸收功率，读数为负表示被测电路发出功率。



\section{主要仪器设备}
\begin{enumerate}[label=\arabic*., leftmargin=2em]
    \item 实验电路板。
    \item 电工电子综合实验台。
    \item 数字式万用表。

\end{enumerate}

\section{操作方法和实验步骤}

\subsection*{1. 交流电压、电流及功率的测量}
用实验电路板上的镇流器替代线圈，白炽灯替代电阻。按图5.4-5接好实验电路，调节交流电源电压 $U$ 至 220V，测量不接电容时 $R_L$ 支路的电流 $I_{RL}$、电源电压 $U$、镇流器（线圈）两端电压 $U_L$、电阻性负载（白炽灯）两端电压 $U_R$ 及电路功率 $P$、功率因数 $\lambda$，记入表1。

\subsection*{2. 并联电容对功率因数的影响}
测量并联不同电容量时的总电流 $I$ 和各支路电流 $I_{RL}$、$I_C$ 及电路功率 $P$、功率因数 $\lambda$，记入表2。注意并联电容 $C$ 要涵盖电路性质是感性和容性两种情况。


\section{实验数据记录和处理}
\subsection{数据记录}

% 表5.4-1 交流电压、电流及功率的测量
\begin{table}[H]
    \centering
    \small
    \begin{tabular}{|c|c|c|c|c|c|c|}
        \hline
        \multicolumn{6}{|c|}{测量値} & 计算値 \\
        \hline
        $U/\mathrm{V}$ & $U_L/\mathrm{V}$ & $U_R/\mathrm{V}$ & $I_{RL}/\mathrm{A}$ & $P/\mathrm{W}$ & $\lambda$ & $\cos\varphi_{RL}$ \\
        \hline
         226.7& 166.8&124.9 &0.245 &38.5 &0.697 & 0.698\\
        \hline
    \end{tabular}
    \caption{交流电压、电流及功率的测量}
    \label{tab:ac_measure}
\end{table}

% 表5.4-2 并联电容对功率因数的影响
\begin{table}[H]
    \centering
    \small
    \begin{tabular}{|c|c|c|c|c|c|c|c|}
        \hline
        \multirow{2}{*}{并联电容 $C/\mu\mathrm{F}$} & \multicolumn{5}{c|}{测量値} & 计算値 & \multirow{2}{*}{电路性质} \\
        \cline{2-7}
         & $I/\mathrm{A}$ & $I_C/\mathrm{A}$ & $I_{RL}/\mathrm{A}$ & $P/\mathrm{W}$ & $\lambda$ & $\cos\varphi$ & \\
        \hline
        0.5 &0.221 &0.0392 &0.246 &38.7 &0.771 &0.772 & 感性\\
        \hline
        1.0 &0.204 &0.0778 &0.245 &38.6 &0.838 & 0.837& 感性\\
        \hline
        1.5 &0.193 &0.1178 &0.245 &38.4 &0.905 &0.904 & 感性\\
        \hline
        2.0 &0.188 &0.1470 &0.244 &38.1 &0.934 &0.934 & 感性\\
        \hline
        2.5 &0.194 &0.1765 &0.244 &38.0 & 0.921& 0.920& 容性\\
        \hline
        3.0 &0.184 &0.220 &0.244 &38.0 &0.881 &0.882 &容性 \\
        \hline
        3.5 &0.204 &0.258 & 0.244&38.3 &0.799 &0.800 & 容性\\
        \hline
        4.0 & 0.225&0.294 &0.244 &38.2 &0.700 &0.700 & 容性\\
        \hline
    \end{tabular}
    \caption{并联电容对功率因数的影响}
    \label{tab:capacitor_pf}
\end{table}


\subsection{数据处理}

对表2、表3、表4中各组测量数据分别进行线性拟合，拟合方程均为 $U_{\mathrm{AB}} = k \cdot I_R + b$的形式。

\section*{六、实验结果与分析}

\section*{七、讨论、心得}

\end{document}



